%% ----------------------------------------------------------------------------
%%  §1.7  Funzioni lineari. Funzione valore assoluto
%% ----------------------------------------------------------------------------
\section{Funzioni lineari. Funzione valore assoluto}
\label{sec:lineari}

\begin{definizione}[Funzione lineare]\index{funzione!lineare}\index{coefficiente angolare}
Si chiama \emph{funzione lineare} (o \emph{funzione affine}) una funzione del tipo
\begin{equation}\label{eq:lineare}
  y=mx+q,
\end{equation}
dove $m,q$ sono numeri reali fissati. Il grafico di una tale funzione è una retta, di cui
il parametro $m$ è detto \emph{coefficiente angolare} e $q$ \emph{ordinata all'origine}
(è il valore $f(0)$).
\end{definizione}

\begin{proposizione}\label{prop:lineare-monotona}
Ogni funzione lineare $f(x)=mx+q$ è monotòna su $\R$; più precisamente è strettamente
crescente se $m>0$, strettamente decrescente se $m<0$ e costante se $m=0$. Se $m\neq0$
essa è invertibile da $\R$ in $\R$, con inversa
\begin{equation}\label{eq:inversa-lineare}
  f^{-1}(y)=\frac{y-q}{m}.
\end{equation}
\end{proposizione}

\begin{schema}
\Ipotesi{$f(x)=mx+q$, con $m,q\in\R$.}
\Tesi{$m>0\implica f$ strettamente crescente; $m<0\implica f$ strettamente decrescente;
$m=0\implica f$ costante; se $m\neq0$, $f$ è invertibile e $f^{-1}(y)=(y-q)/m$.}
\end{schema}

\begin{strategia}
È un'applicazione diretta della \cref{prop:ordine-prodotto}: moltiplicare la
disuguaglianza $x_{1}<x_{2}$ per $m$ la conserva se $m>0$ e la rovescia se $m<0$;
sommare $q$ non cambia nulla. Per l'inversa basta risolvere l'equazione $y=mx+q$
rispetto a $x$, cosa possibile se e solo se $m\neq0$.
\end{strategia}

\begin{dimostrazione}
Consideriamo $x_{1}<x_{2}$; allora
\[
  f(x_{1})=mx_{1}+q, \qquad f(x_{2})=mx_{2}+q .
\]
Se il coefficiente angolare $m$ è positivo, per la \cref{prop:ordine-prodotto}(i) da
$x_{1}<x_{2}$ segue $mx_{1}<mx_{2}$ (la disuguaglianza resta stretta perché $m\neq0$ e
la \cref{prop:semplificazione-prodotto} esclude l'uguaglianza) e quindi, sommando $q$
(assioma \ref{ass:ordine-somma}), $f(x_{1})<f(x_{2})$: $f$ è strettamente crescente. Se
$m$ è negativo, per la \cref{prop:ordine-prodotto}(ii) da $x_{1}<x_{2}$ segue
$mx_{1}>mx_{2}$ e quindi $f(x_{1})>f(x_{2})$: $f$ è strettamente decrescente. Infine, se
$m=0$, risulta $f(x)=q$ per ogni $x$: la funzione è costante, dunque
contemporaneamente crescente e decrescente, e il suo grafico è una retta parallela
all'asse $x$.

Se $m\neq0$, $f$ è strettamente monotòna; è anche suriettiva su $\R$, perché per ogni
$y\in\R$ l'equazione $mx+q=y$ ha la soluzione $x=(y-q)/m$. Per il
\cref{teo:monotona-invertibile} $f$ è invertibile, e la \eqref{eq:inversa-lineare} è
di verifica immediata: $f\bigl((y-q)/m\bigr)=m\cdot(y-q)/m+q=y$. Al contrario, se $m=0$
la funzione costante $f(x)=q$ non stabilisce una corrispondenza biunivoca tra $\R$ e
l'insieme $\{q\}$, e non è invertibile.
\end{dimostrazione}

Dato che per due punti distinti del piano passa una e una sola retta, per disegnare il
grafico di una funzione lineare è sufficiente calcolare l'ordinata $y$ in
\eqref{eq:lineare} in corrispondenza a due valori distinti della variabile $x$. Ad
esempio, nel caso della funzione lineare $y=3x-1$, a $x=0$ corrisponde $y=-1$ e a $x=1$
corrisponde $y=2$; il grafico è la retta passante per i punti $(0,-1)$ e $(1,2)$
(\cref{fig:rette}).

\begin{figure}[htb]
\centering
\begin{tikzpicture}
\begin{axis}[width=5.8cm, height=5.2cm, xmin=-1.5, xmax=2.5, ymin=-2.5, ymax=4,
  xtick={1}, ytick={-1,2}, title={\footnotesize $y=3x-1$}, title style={at={(0.3,0.9)}}]
  \addplot[curva, domain=-0.5:1.6] {3*x-1};
  \addplot[tratteggio] coordinates {(1,0) (1,2) (0,2)};
  \addplot[only marks, mark=*, mark size=1.4pt, rame] coordinates {(0,-1) (1,2)};
\end{axis}
\end{tikzpicture}\hfill
\begin{tikzpicture}
\begin{axis}[width=5.8cm, height=5.2cm, xmin=-1, xmax=3.5, ymin=-1.5, ymax=4,
  xtick={1,2,3}, ytick={1,2,3}, title={\footnotesize $y=3-2x$}, title style={at={(0.75,0.9)}}]
  \addplot[curva, domain=-0.5:2.2] {3-2*x};
  \addplot[only marks, mark=*, mark size=1.4pt, rame] coordinates {(0,3) (1.5,0)};
\end{axis}
\end{tikzpicture}
\caption{Grafici di due funzioni lineari: $y=3x-1$ (coefficiente angolare $m=3>0$,
strettamente crescente) e $y=3-2x$ ($m=-2<0$, strettamente decrescente). In rame i
due punti usati per tracciare ciascuna retta.}
\label{fig:rette}
\end{figure}

\begin{idea}[Il coefficiente angolare come «velocità»]
Il coefficiente angolare $m$ misura di quanto cresce $y$ quando $x$ cresce di $1$:
$f(x+1)-f(x)=m$. Se $x$ è il tempo e $y$ la posizione di un carrello lungo una guida,
$m$ è la velocità: positiva se il carrello avanza, negativa se arretra, nulla se è fermo.
La \cref{prop:lineare-monotona} dice allora qualcosa di ovvio dal punto di vista
fisico: un carrello con velocità costante positiva non torna mai indietro. Il calcolo
differenziale, nei capitoli successivi, estenderà questa idea a funzioni qualsiasi,
definendo una «velocità istantanea» --- la derivata --- che varia da punto a punto.
\end{idea}

\subsection{Il valore assoluto}

\begin{definizione}[Valore assoluto]\index{valore assoluto}\index{modulo}
Il \emph{valore assoluto} (o \emph{modulo}) di $x\in\R$, indicato con il simbolo
$\abs{x}$, è definito da
\begin{equation}\label{eq:modulo}
  \abs{x}=\begin{cases} x & \text{se } x\ge0,\\ -x & \text{se } x<0. \end{cases}
\end{equation}
\end{definizione}

\begin{marginfigure}
\centering
\begin{tikzpicture}
\begin{axis}[width=5cm, height=4cm, xmin=-2.5, xmax=2.5, ymin=-0.5, ymax=2.5,
  xtick={-1,1}, ytick={1}, title={\footnotesize $y=\abs{x}$}, title style={at={(0.5,0.9)}}]
  \addplot[curva, domain=-2.3:0] {-x};
  \addplot[curva, domain=0:2.3] {x};
\end{axis}
\end{tikzpicture}
\caption{Il grafico della funzione valore assoluto è l'unione della semiretta $y=x$,
per $x\ge0$, e della semiretta $y=-x$, per $x<0$.}
\label{fig:modulo}
\end{marginfigure}

Il grafico della funzione valore assoluto $f(x)=\abs{x}$ è composto da due semirette
uscenti dall'origine, di equazione rispettivamente $y=x$ e $y=-x$, come nella
\cref{fig:modulo}. Geometricamente, $\abs{x}$ è la \emph{distanza} del punto $x$
dall'origine sulla retta reale, e $\abs{x-y}$ è la distanza tra i punti $x$ e $y$: è
questa interpretazione che rende il valore assoluto onnipresente in Analisi, dove
«vicino» significa sempre «a distanza piccola».

\begin{proposizione}[Proprietà elementari del valore assoluto]\label{prop:modulo-proprieta}
Per ogni $x,x_{1},x_{2}\in\R$:
\begin{align}
  &\abs{x}\ge0; \label{eq:mod1}\\
  &\abs{x}=0 \ \sse\ x=0; \label{eq:mod2}\\
  &\abs{-x}=\abs{x}; \label{eq:mod3}\\
  &\abs{x_{1}\cdot x_{2}}=\abs{x_{1}}\cdot\abs{x_{2}}; \label{eq:mod4}\\
  &\abs{x_{1}/x_{2}}=\abs{x_{1}}/\abs{x_{2}} \quad (x_{2}\neq0). \label{eq:mod5}
\end{align}
\end{proposizione}

\begin{schema}
\Ipotesi{$x,x_{1},x_{2}\in\R$; nella \eqref{eq:mod5}, $x_{2}\neq0$.}
\Tesi{Le relazioni \eqref{eq:mod1}--\eqref{eq:mod5}.}
\end{schema}

\begin{strategia}
Il valore assoluto è definito \emph{per casi}, e le sue proprietà si dimostrano quindi
\emph{distinguendo i casi}: $x\ge0$ e $x<0$ (e, per le proprietà con due variabili, le
quattro combinazioni di segno). È un metodo noioso ma infallibile; conviene farlo una
volta per tutte e poi non pensarci più.
\end{strategia}

\begin{dimostrazione}
Dimostriamo ad esempio la \eqref{eq:mod1}: se $x\ge0$ allora $\abs{x}=x\ge0$; se invece
$x<0$, allora $\abs{x}=-x$ e, per la \cref{prop:segno-opposto}, $-x>0$; in definitiva
$\abs{x}\ge0$ per ogni $x\in\R$. La \eqref{eq:mod2} è immediata dalla definizione:
$\abs{x}=0$ significa $x=0$ (caso $x\ge0$) oppure $-x=0$ (caso $x<0$, impossibile).
Proviamo la \eqref{eq:mod3}: se $x>0$ allora $\abs{x}=x$ e $\abs{-x}=-(-x)=x$, dato che
$-x<0$; se $x<0$ risulta $\abs{x}=-x$ e $\abs{-x}=-x$, essendo $-x>0$; infine, se $x=0$,
si ha $-0=0$ e pertanto $\abs{-0}=\abs{0}$. La \eqref{eq:mod4} è diretta conseguenza
della regola dei segni (\cref{prop:segni}): ad esempio, se $x_{1}<0$ e $x_{2}\ge0$, il
prodotto $x_{1}x_{2}$ è $\le0$, dunque
$\abs{x_{1}x_{2}}=-(x_{1}x_{2})=(-x_{1})x_{2}=\abs{x_{1}}\abs{x_{2}}$; gli altri casi
sono analoghi. Infine la \eqref{eq:mod5} segue dalla \eqref{eq:mod4}:
$\abs{x_{1}}=\abs{(x_{1}/x_{2})\,x_{2}}=\abs{x_{1}/x_{2}}\,\abs{x_{2}}$, e basta dividere
per $\abs{x_{2}}$, che è non nullo per la \eqref{eq:mod2}.
\end{dimostrazione}

Esaminiamo ora due proprietà la cui interpretazione geometrica è schematizzata nella
\cref{fig:modulo-r}: esse permettono di tradurre una disuguaglianza con il valore
assoluto in una \emph{doppia} disuguaglianza senza valore assoluto, ed è così che le
useremo sistematicamente nello studio dei limiti.

\begin{proposizione}\label{prop:modulo-r}
Per ogni numero reale $r\ge0$ valgono le equivalenze:
\begin{align}
  \abs{x}\le r &\ \sse\ -r\le x\le r; \label{eq:mod-le}\\
  \abs{x}<r    &\ \sse\ -r<x<r. \label{eq:mod-lt}
\end{align}
\end{proposizione}

\begin{schema}
\Ipotesi{$r\ge0$, $x\in\R$.}
\Tesi{$\abs{x}\le r$ se e solo se $x\le r$ \emph{e} $x\ge-r$ (e analogamente con le
disuguaglianze strette).}
\end{schema}

\begin{strategia}
Ancora una distinzione di casi: la relazione $\abs{x}\le r$ significa «$x\le r$ se
$x\ge0$» oppure «$-x\le r$ se $x<0$». Ciascun caso dà un pezzo dell'intervallo
$[-r,r]$: il primo la metà destra $[0,r]$, il secondo la metà sinistra $[-r,0)$.
Riunendo i due pezzi si ottiene tutto l'intervallo.
\end{strategia}

\begin{dimostrazione}
La verifica delle due equivalenze procede allo stesso modo; perciò ci limitiamo a
dimostrare la \eqref{eq:mod-le}. In base alla definizione \eqref{eq:modulo} del valore
assoluto, la relazione $\abs{x}\le r$ equivale a uno dei due sistemi
\begin{equation}\label{eq:mod-sistemi}
  \begin{cases} x\ge0\\ x\le r\end{cases} \qquad\text{oppure}\qquad
  \begin{cases} x<0\\ -x\le r.\end{cases}
\end{equation}
Il primo sistema si riscrive nella forma $0\le x\le r$, mentre il secondo equivale (per
la \cref{prop:ordine-prodotto}(ii), moltiplicando per $-1$) a $-r\le x<0$. Unendo i due
risultati si trova infine $-r\le x\le r$. Viceversa, se $-r\le x\le r$, allora
$x\le r$ e $-x\le r$; poiché $\abs{x}$ è uguale a $x$ oppure a $-x$, in ogni caso
$\abs{x}\le r$.
\end{dimostrazione}

\begin{figure}[htb]
\centering
\begin{tikzpicture}
\begin{axis}[width=8cm, height=5cm, xmin=-2.8, xmax=2.8, ymin=-0.4, ymax=2.6,
  xtick={-1.6,1.6}, xticklabels={$-r$,$r$}, ytick={1.6}, yticklabels={$r$}]
  \addplot[curva, domain=-2.5:0] {-x};
  \addplot[curva, domain=0:2.5] {x};
  \addplot[tratteggio] coordinates {(-1.6,0) (-1.6,1.6) (1.6,1.6) (1.6,0)};
  \addplot[rame, line width=2.2pt] coordinates {(-1.6,0) (1.6,0)};
  \addplot[verde, line width=2.2pt] coordinates {(0,0) (0,1.6)};
  \node[font=\scriptsize, rame, below] at (axis cs:0.85,-0.02) {$-r\le x\le r$};
  \node[font=\scriptsize, verde, right] at (axis cs:0.05,0.9) {$y\le r$};
  \node[font=\scriptsize, blu, anchor=east] at (axis cs:2.05,2.35) {$y=\abs{x}$};
\end{axis}
\end{tikzpicture}
\caption{Interpretazione geometrica della \cref{prop:modulo-r}: i punti del grafico di
$y=\abs{x}$ con ordinata $y\le r$ (segmento verde) sono esattamente quelli con ascissa
nell'intervallo $[-r,r]$ (segmento color rame).}
\label{fig:modulo-r}
\end{figure}

La seguente proprietà del valore assoluto, detta \emph{disuguaglianza triangolare}, è
di grande importanza. La spiegazione intuitiva è semplice: a secondo membro della
\eqref{eq:triangolare} compare la somma di due numeri positivi o nulli (i moduli dei
numeri $x_{1}$ e $x_{2}$), mentre a primo membro compare il modulo della \emph{somma
algebrica} di $x_{1}$ e $x_{2}$; se i segni di $x_{1}$ e $x_{2}$ sono discordi, il
primo membro è minore del secondo, mentre se i segni sono concordi i due membri sono
uguali.

\begin{teorema}[Disuguaglianza triangolare]\label{teo:triangolare}\index{disuguaglianza triangolare}
Per ogni coppia di numeri reali $x_{1},x_{2}$ vale la disuguaglianza
\begin{equation}\label{eq:triangolare}
  \abs{x_{1}+x_{2}}\le\abs{x_{1}}+\abs{x_{2}}.
\end{equation}
\end{teorema}

\begin{schema}
\Ipotesi{$x_{1},x_{2}\in\R$.}
\Tesi{$\abs{x_{1}+x_{2}}\le\abs{x_{1}}+\abs{x_{2}}$.}
\end{schema}

\begin{strategia}
Potremmo distinguere i quattro casi di segno, ma c'è una strada più elegante che
sfrutta la \cref{prop:modulo-r} \emph{due volte}: una volta «in uscita», per
trasformare le ovvie relazioni $\abs{x_{i}}\le\abs{x_{i}}$ nelle doppie disuguaglianze
$-\abs{x_{i}}\le x_{i}\le\abs{x_{i}}$; poi sommiamo membro a membro; infine usiamo la
\cref{prop:modulo-r} «in entrata», per ricompattare la doppia disuguaglianza ottenuta
in una disuguaglianza con il modulo. È un esempio tipico di come una
\emph{equivalenza} si possa usare nei due versi.
\end{strategia}

\begin{dimostrazione}
Per ogni numero reale $x$ la relazione $\abs{x}\le\abs{x}$ è ovvia (anzi, vale con il
segno di uguaglianza); indicando con $r$ il secondo membro di tale ovvia relazione, per
l'equivalenza \eqref{eq:mod-le} abbiamo anche $-\abs{x}\le x\le\abs{x}$; in particolare,
per $x=x_{1}$ e $x=x_{2}$:
\begin{equation}\label{eq:triang-passo1}
  -\abs{x_{1}}\le x_{1}\le\abs{x_{1}}, \qquad -\abs{x_{2}}\le x_{2}\le\abs{x_{2}}.
\end{equation}
Sommando membro a membro (è lecito: per \ref{ass:ordine-somma} e la proprietà
transitiva, da $a\le b$ e $c\le d$ segue $a+c\le b+c\le b+d$):
\begin{equation}\label{eq:triang-passo2}
  -\bigl(\abs{x_{1}}+\abs{x_{2}}\bigr)\le x_{1}+x_{2}\le\abs{x_{1}}+\abs{x_{2}}.
\end{equation}
La conclusione segue da una nuova applicazione dell'equivalenza \eqref{eq:mod-le}, con
$r=\abs{x_{1}}+\abs{x_{2}}$ e $x=x_{1}+x_{2}$.
\end{dimostrazione}

\begin{notastorica}[Perché «triangolare»?]
Il nome viene dalla geometria: in un triangolo, ogni lato è minore della somma degli
altri due (proposizione I.20 degli \emph{Elementi} di Euclide, che il filosofo Proclo
riferisce fosse derisa dagli epicurei perché «nota anche a un asino», che per
raggiungere il fieno va dritto e non lungo i due lati). Interpretando $x_{1}$ e $x_{2}$
come vettori nel piano, $\abs{x_{1}+x_{2}}\le\abs{x_{1}}+\abs{x_{2}}$ è esattamente
quella proposizione; sulla retta i «triangoli» sono degeneri, ma il nome è rimasto. La
disuguaglianza triangolare è la proprietà fondamentale di ogni nozione di
\emph{distanza}: è uno degli assiomi con cui Maurice Fréchet, nel 1906, introdusse gli
spazi che oggi chiamiamo \emph{metrici}.
\end{notastorica}

\begin{corollario}\label{cor:triangolare-inversa}
Per ogni $x_{1},x_{2}\in\R$ risulta
\begin{equation}\label{eq:triangolare-inversa}
  \bigl|\abs{x_{1}}-\abs{x_{2}}\bigr|\le\abs{x_{1}-x_{2}}.
\end{equation}
\end{corollario}

\begin{schema}
\Ipotesi{$x_{1},x_{2}\in\R$.}
\Tesi{$\bigl|\abs{x_{1}}-\abs{x_{2}}\bigr|\le\abs{x_{1}-x_{2}}$.}
\end{schema}

\begin{strategia}
Vogliamo una stima di $\abs{x_{1}}-\abs{x_{2}}$: il trucco è scrivere $x_{1}$ come
$(x_{1}-x_{2})+x_{2}$ e applicare la disuguaglianza triangolare a questa somma. Si
ottiene $\abs{x_{1}}-\abs{x_{2}}\le\abs{x_{1}-x_{2}}$; scambiando i ruoli si ottiene
la stima con il segno opposto, e la \cref{prop:modulo-r} ricompone le due.
\end{strategia}

\begin{dimostrazione}
Per il \cref{teo:triangolare}, $\abs{x_{1}}=\abs{(x_{1}-x_{2})+x_{2}}\le\abs{x_{1}-x_{2}}+\abs{x_{2}}$,
da cui $\abs{x_{1}}-\abs{x_{2}}\le\abs{x_{1}-x_{2}}$. Scambiando $x_{1}$ con $x_{2}$ e
usando la \eqref{eq:mod3}: $\abs{x_{2}}-\abs{x_{1}}\le\abs{x_{2}-x_{1}}=\abs{x_{1}-x_{2}}$,
cioè $-\abs{x_{1}-x_{2}}\le\abs{x_{1}}-\abs{x_{2}}$. Le due stime, per la
\eqref{eq:mod-le} con $r=\abs{x_{1}-x_{2}}$, danno la tesi.
\end{dimostrazione}

\begin{python}[Verifica sperimentale della disuguaglianza triangolare]
# Un esperimento numerico non dimostra nulla, ma aiuta a "vedere"
# quando la disuguaglianza è stretta e quando è un'uguaglianza.
import random
random.seed(1)
for _ in range(6):
    x1, x2 = random.uniform(-5, 5), random.uniform(-5, 5)
    sinistra, destra = abs(x1 + x2), abs(x1) + abs(x2)
    segni = "concordi" if x1 * x2 >= 0 else "discordi"
    rel = "=" if abs(sinistra - destra) < 1e-12 else "<"
    print(f"x1={x1:+.2f} x2={x2:+.2f} ({segni}): "
          f"|x1+x2|={sinistra:.2f} {rel} |x1|+|x2|={destra:.2f}")
\end{python}
\begin{risultato}
x1=-3.66 x2=+3.47 (discordi): |x1+x2|=0.18 < |x1|+|x2|=7.13
x1=+2.64 x2=-2.45 (discordi): |x1+x2|=0.19 < |x1|+|x2|=5.09
x1=-0.05 x2=-0.51 (concordi): |x1+x2|=0.55 = |x1|+|x2|=0.55
x1=+1.52 x2=+2.89 (concordi): |x1+x2|=4.40 = |x1|+|x2|=4.40
x1=-4.06 x2=-4.72 (concordi): |x1+x2|=8.78 = |x1|+|x2|=8.78
x1=+3.36 x2=-0.67 (discordi): |x1+x2|=2.69 < |x1|+|x2|=4.03
\end{risultato}

\begin{esercizio}
Risolvere la disequazione $\abs{2x-3}<5$.
\begin{soluzione}
Per la \eqref{eq:mod-lt} con $r=5$: $-5<2x-3<5$, cioè $-2<2x<8$, cioè $-1<x<4$.
\emph{Scelta logica:} la \cref{prop:modulo-r} trasforma la disequazione con il modulo
in una doppia disequazione senza modulo, ed evita la discussione dei due casi
$2x-3\ge0$ e $2x-3<0$.
\end{soluzione}
\end{esercizio}

\begin{esercizio}
Dimostrare che per ogni $x,y\in\R$ si ha $\max\{x,y\}=\dfrac{x+y+\abs{x-y}}{2}$ e
$\min\{x,y\}=\dfrac{x+y-\abs{x-y}}{2}$.
\begin{soluzione}
I secondi membri non cambiano scambiando $x$ con $y$, perché $\abs{x-y}=\abs{y-x}$
(\eqref{eq:mod3}); lo stesso vale per $\max\{x,y\}$ e $\min\{x,y\}$. Possiamo quindi
supporre $x\ge y$. Allora $x-y\ge0$, $\abs{x-y}=x-y$, e
\[
\begin{aligned}
  \frac{x+y+\abs{x-y}}{2}&=\frac{x+y+(x-y)}{2}=x=\max\{x,y\},\\
  \frac{x+y-\abs{x-y}}{2}&=\frac{x+y-(x-y)}{2}=y=\min\{x,y\}.
\end{aligned}
\]
\emph{Scelta logica:} il modulo è definito per casi, e il caso da trattare dipende dal
segno del suo argomento $x-y$; la simmetria delle formule rispetto allo scambio di $x$
con $y$ permette di trattarne uno solo, «senza perdita di generalità».
\end{soluzione}
\end{esercizio}
